\documentclass{article}
\usepackage{mathtools} 
\usepackage{fontspec}
\usepackage[UTF8]{ctex}
\usepackage{amsthm}
\usepackage{mdframed}
\usepackage{xcolor}
\usepackage{amssymb}
\usepackage{amsmath}


% 定义新的带灰色背景的说明环境 zremark
\newmdtheoremenv[
  backgroundcolor=gray!10,
  % 边框与背景一致，边框线会消失
  linecolor=gray!10
]{zremark}{说明}

% 通用矩阵命令: \flexmatrix{矩阵名}{元素符号}{行数}{列数}
\newcommand{\flexmatrix}[4]{
  \[
  #1 = \begin{pmatrix}
    #2_{11}     & #2_{12}     & \cdots & #2_{1#4}   \\
    #2_{21}     & #2_{22}     & \cdots & #2_{2#4}   \\
    \vdots      & \vdots      & \ddots & \vdots     \\
    #2_{#31}    & #2_{#32}    & \cdots & #2_{#3#4}
  \end{pmatrix}
  \]
}

% 简化版命令（默认矩阵名为A，元素符号为a）: \quickmatrix{行数}{列数}
\newcommand{\quickmatrix}[2]{\flexmatrix{A}{a}{#1}{#2}}

\begin{document}
\title{习题 17.5}
\author{张志聪}
\maketitle

\section*{5}
以$D$是开区域为例。

(i) 如果$D$是凸域，由对任意$P_0(x_0, y_0), P_1(x_1, y_1) \in D$，
由定理17.8（中值定理）可得，存在$0 < \theta < 1$，使得
\begin{align*}
  f(x_0, y_0) - f(x_1, y_1)
   & = f_x(x_0 + \theta(x_1 - x_0), y_0 + \theta(y_1 - y_0)) + f_y(x_0 + \theta(x_1 - x_0), y_0 + \theta(y_1 - y_0)) \\
   & = 0 + 0                                                                                                         \\
   & = 0
\end{align*}
所以$f(x_0, y_0) = f(x_1, y_1)$，故$f$在$D$上是常值函数。

(ii) 如果$D$不是凸域，现在任取$(x_0, y_0) \in D$，令
\begin{align*}
  U = \{(x, y) \in D : f(x, y) = f(x_0, y_0) \}
\end{align*}
显然$U$不是空集。

我们证明$U$是开集。任取$(a_1, a_2) \in U \subseteq D$，由于$D$是开集，
存在以$(a_1, a_2)$为原形的球（邻域）$B_1 \subseteq D$，球$B_1$是凸域，
由(i)可知，$f$在$B_1$上是常值函数：对任意$(x, y) \in B_1$，有
\begin{align*}
  f(x, y) = f(a_1, a_2) = f(x_0, y_0)
\end{align*}
$B_1 \subseteq U$，$(a_1, a_2)$是$U$的内点，所以$U$是开集。

同样可证$D \cap U^c$是一个开集，这是因为任取一点$(b_1, b_2) \subseteq D \cap U^c$，
存在以$(b_1, b_2)$为原形的球（邻域）$B_2 \subseteq D$，球$B_2$是凸域，
由(i)可知，$f$在$B_2$上是常值函数：对任意$(x, y) \in B_2$，有
\begin{align*}
  f(x, y) = f(b_1, b_2) \neq f(x_0, y_0)
\end{align*}
所以$B_2$中的点不在$U$中，
这表明$B_2 \subseteq D \cap U^c$，
$(b_1, b_2)$是$D \cap U^c$的内点，所以$B_2$是开集。

我们有
\begin{align*}
 D = U \cup (D \cap U^c) 
\end{align*}
由于$D$是连通的开集和$U \neq \varnothing$，则必有
$D \cap U^c = \varnothing$（这里使用了连通性的等价定义：$X$是连通的，
当且仅当$X$中不存在两个不相交的非空开集$V, W$使得$V \cup W = X$），
所以$D = U$，故$f$在$D$上是常值函数。




\end{document}